% Book Two Pass 3 cumulative project, additive after the original final chapter.
\clearpage
\section*{Cumulative project: the Numerical Geometry Atlas}
\addcontentsline{toc}{section}{Cumulative project: the Numerical Geometry Atlas}

\noindent A mathematical idea becomes an engineering capability when we can
derive its governing equations, implement the algorithm, recognize its
assumption failures, and explain the numerical answer without exaggeration.
The Numerical Geometry Atlas is the capstone built from the source-led
laboratories of this book. Each scientific plate has a program that generated
its data, mathematical reference where available, units, failure conditions,
and an independent reader reconstruction exercise. The seven figures are
not copied from the original research papers.

\paragraph{The source code and evidence chain.}
The \texttt{tools/labs} directory contains the original small numerical
implementations used in this course. The geometry laboratory checks
orientation signs and performs reproducible two-dimensional RANSAC fitting.
The numerical laboratory assembles triangle stiffness and lumped-mass
matrices, computes a heat-method distance field, and independently fits
a nonlinear model under Huber loss. The sampling laboratory uses normalized
proposal densities to estimate an integral with exact analytical value.
A C++23 companion checks sign certification, discrete energy, numerical
conditioning and safe handling of uncertain answers.

The renderer \texttt{book02\_visual\_atlas.py} calls the actual mathematical
implementations and produces three visual forms: PDF vectors for print, SVG
for scalable screens and digital publishing, and PNG for previews. It saves
every figure's underlying numerical rows as CSV, with a JSON manifest
recording units, independent references, seeded experiments, library versions
and limits of inference. The companion HTML atlas is generated from exactly
those exported rows, with all required data embedded. It works offline:
a control can rotate a known computed surface or change the scalar field
shown, but does not invent a new solve or silently recalculate a research
result that was never produced by the source implementation.

\paragraph{Checkpoint I: classify geometric orientation reliably.}
For planar coordinates \(a,b,c\), derive the oriented-area predicate
\[
D=(a_x-c_x)(b_y-c_y)-(a_y-c_y)(b_x-c_x).
\]
Test a clockwise triangle, a counterclockwise triangle, exact collinearity
and a near-collinear binary64 input. Explain when the floating-point filter
can certify a sign and when the rational exact-sign fallback is needed.
An uncertain C++23 answer must not silently be classified as zero.
Reversing vertex order must reverse the sign and a valid translation of
all three vertices must preserve it. Identify a case in which intermediate
product overflow makes the naive determinant unusable.

\paragraph{Checkpoint II: consensus under outliers.}
Create a synthetic scene containing 80 near-line samples and 30 structured
outliers. Implement minimal two-point line fitting, threshold-based
consensus and orthogonal total-least-squares refinement. Compare the
robust result with the all-data TLS fit, and disclose the threshold and seed.
The offline viewer changes consensus membership against a fixed fitted
line; an actual refit requires running the numerical source program again.
Compare false acceptance, false rejection, parameter error, and stability
over separately seeded experiments. No one plot is evidence of universal
RANSAC superiority.

\paragraph{Checkpoint III: discretize a continuous geometric field.}
For linear basis functions \(\phi_i\) on a triangle \(T\), derive
\[
u(x)=\sum_i u_i\phi_i(x),\qquad
(K_T)_{ij}=A_T\nabla\phi_i\cdot\nabla\phi_j,\qquad
E_T(u)=\tfrac12u^TK_Tu.
\]
Assemble the sparse positive-semidefinite stiffness matrix and a diagonal
lumped mass matrix. Check its symmetry, the constant-vector nullspace,
total mass and nonnegative energy. A discretization may reduce energy yet
violate a maximum principle; those are different mathematical properties.
Document matrix signs, units, boundary assumptions and degeneracy handling.

\paragraph{Checkpoint IV: recover geodesics from heat.}
Take a backward implicit heat step,
\[
(M+tK)u=Mu_0,
\]
normalize the negative triangle-wise temperature gradient, then solve
an anchored Poisson equation for distance. Compare the recovered distances
against the Euclidean reference on a flat convex square. Save the signed
field error and absolute error separately. Display computed heat, recovered
distance and signed error side by side. The additional 3D picture lifts
distance as a vertical scalar and must explicitly say that the height
does not represent physical terrain.

Repeat over several grid resolutions on the same physical domain.
Record vertex and triangle counts, spacing, mean and median absolute
errors, and maximum absolute error. Since the heat timestep is tied
to mean edge length squared, its scaling is part of the experiment.
Explain any nonmonotonic result rather than deleting inconvenient
points. Move the source and change the triangulation to investigate
the limitations of the baseline.

\paragraph{Checkpoint V: robust nonlinear fitting and identity.}
Generate samples of \(y=a e^{bx}+c\) with deterministic perturbations
and one known gross outlier. Derive the residual Jacobian
\[
J_i=\begin{pmatrix}e^{bx_i}&a x_i e^{bx_i}&1\end{pmatrix}.
\]
Compare ordinary least squares with a damped Huber method and the
SciPy robust-loss reference. Record the final parameters, value,
gradient norm, accepted steps and termination status. Include all
observations in the visual and residual statistics. The estimator
with the smallest squared residual is not automatically the method
that best recovers the generating parameters, because the objective
functions and contamination assumptions differ.

\paragraph{Checkpoint VI: a sampling study that measures uncertainty.}
For the exact integral \(\int_0^1 e^x\,dx=e-1\), compare uniform,
tilted-density and balance-heuristic estimators at equal numbers of
integrand evaluations. Repeat with distinct pseudorandom streams
and visualize the distributions of signed errors and their sample
variance. The purpose is to understand support, measure, weighting,
bias and cost, not to claim that one sampler dominates for every
integrand. Then design a two-peak function for which the proposals
have distinct strengths, and compare over an independently justified
sample and replication budget.

\paragraph{Scientific figure acceptance.}
Every exported plate must have readable axes, correct units, a
declared algorithm and seed, a reference formula or validation
baseline, and an explicit distinction between synthetic data,
analytical truth and solver output. A heatmap colour scale must
say whether it represents signed error or magnitude. A 3D surface
must identify its vertical variable. Never invent hardware speedups,
paper-reproduction figures or uncertainty intervals for an attractive
presentation. When measuring runtime later, record machine, compiler,
compiler flags, software version, workload, independent repeats,
hardware state and variance.

A checkpoint is accepted only if its tests pass, its output is
reproducible, invalid inputs fail explicitly, and the numerical
results in the CSV match the plotted data. The offline viewer must
work without a network connection and use the exported tables,
not a separate visual approximation. The reader must be able
to reconstruct at least one result by hand and one by code.

\paragraph{Beyond the initial atlas.}
Extend the heat method to alternative domains and boundaries.
Compare adaptive floating expansions with the rational sign oracle.
Test RANSAC with competing structures or very low inlier rates.
Move nonlinear fitting from a curve to a geometric pose problem with
explicit gauge freedoms. Every extension should include a new
derivation, adversarial test, independent reference, reconstructed
figure and a discussion of limits. The object is not to assemble
the most colourful demonstration; it is to build numerical software
whose mathematical meaning can be defended.
